%% ma: L12-B01-0034
%% lop: 12
%% chuong: 1
%% bai: 1
%% dang: 12.1.2
%% loai: TN4
%% muc: NB
%% dap_an: "A"
%% nguon: "HK1-DeOn_1_1.pdf (M10, Đề ôn tập học kì 1 Toán 12, Đề số 2), Phần I Câu 3"
%% kien_thuc: "Bài 1 (tính đơn điệu và cực trị)"
%% trang_thai: chinh-thuc
%% ghi_chu: "Hình vẽ lại theo đồ thị y = x³ + 3x² − 2 (cực đại (−2; 2), cực tiểu (0; −2) như hình gốc)."
\begin{ex}
Cho hàm số $y=f(x)$ có đồ thị như hình vẽ:
\begin{center}\begin{tikzpicture}[x=0.9cm,y=0.6cm]
\draw[truc] (-3.4,0)--(1.8,0) node[below] {$x$};
\draw[truc] (0,-3)--(0,4.2) node[left] {$y$};
\node[below right] at (0,0) {$O$};
\draw[phu] (-2,0)--(-2,2)--(0,2);
\node[below] at (-2,0) {$-2$}; \node[right] at (0,2) {$2$}; \node[below right] at (0,-2) {$-2$};
\draw[thick,smooth,samples=100,domain=-3.15:0.95] plot (\x,{\x*\x*\x+3*\x*\x-2});
\end{tikzpicture}\end{center}
Hàm số $y=f(x)$ đồng biến trên khoảng nào trong các khoảng sau đây?
\choice{\True $(-\infty;-2)$}{$(-2;0)$}{$(-2;+\infty)$}{$(-\infty;2)$}
\loigiai{Từ đồ thị, hàm số đồng biến trên $(-\infty;-2)$ và $(0;+\infty)$, nghịch biến trên $(-2;0)$.}
\end{ex}
